% 1-1.5 pages
\section{Introduction} \label{sec:introduction}
Large-scale transformer-based language models have demonstrated remarkable capabilities across a wide range of natural-language tasks, including text generation, summarization, and creative writing \cite{attention}\cite{GPT-3}.
However, these advances rely on extensive computational resources and carefully curated training datasets, making such models costly to train and deploy, therefore often depending on centralized infrastructure \cite{energy}.
These limitations have led to increasing interest in small, efficient models that can operate locally while still supporting domain-specific or personalized applications \cite{smallLMs}.
\\
In this context, our project explores the development of a compact language model designed to generate film reviews from short plot summaries and user input.
Our approach focuses on mimicking user-generated writing styles using a lightweight architecture suitable for deployment on mobile or embedded devices.
To achieve this, we trained the model on a dataset of film synopsis and reviews sourced from the online platform Letterboxd \cite{letterboxd}.
Since these data originate from real users, significant preprocessing was required to sanitize the dataset.
\\
This work therefore focuses both on the limitations of small language models and on the challenging use of user-generated content for training purposes, resulting in the following research questions.


\subsection*{Research Questions} \label{subsec:research-questions}
\begin{enumerate}
    \setlength{\topsep}{0pt}
    \setlength{\itemsep}{0pt}
    \setlength{\parskip}{0pt}
    \item Is it possible for a small LLM to generate concise film reviews when trained on user-written ones?
    \item How must the data be improved to improve the models performance?
    \item Does giving the model a summary of the film improve its abilities to fitting reviews?
    \item What parameters should the model have, specifically how large should it be?
\end{enumerate}


\subsection*{Outline} \label{subsec:outline}
This paper will start by analyzing the dataset and stating how it was improved in section \ref{sec:data}.
In the following section \ref{sec:model}, the used model architectures and metrics will be compared.
Then section \ref{sec:evaluation} explains the evaluation setup, before section \ref{sec:results} shows the results of our experiments, focusing on what changes improved the performance of the model best.
At last, section \ref{sec:conclusion} concludes our findings.